% ============================================================================
% Math4AI Capstone Report — ACM sigconf style
% Sharaf Feyzullayev · Nijat Samadov · Samir Abdullazade
% National AI Center — AI Academy, 2026
% ============================================================================

\documentclass[sigconf,screen,nonacm]{acmart}

% --- ACM metadata cleanup for course-project use ---
\settopmatter{printacmref=false}
\renewcommand\footnotetextcopyrightpermission[1]{}
\pagestyle{plain}
\setcopyright{none}

% --- Project metadata ---
\newcommand{\ProjectTitle}{From Linear Scores to a Single Hidden Layer:\\%
  A Mathematical Study of Simple Learning Systems}
\newcommand{\CourseName}{National AI Center --- AI Academy}
\newcommand{\ProjectType}{Math4AI Final Capstone Project}

% --- Packages and local macros ---
\usepackage[T1]{fontenc}
\usepackage[utf8]{inputenc}
\usepackage{lmodern}
\usepackage{microtype}
\usepackage{amsmath,amssymb,amsfonts,mathtools,bm}
\usepackage{booktabs}
\usepackage{multirow}
\usepackage{array}
\usepackage{graphicx}
\usepackage{subcaption}
\usepackage{xcolor}
\usepackage{enumitem}
\usepackage{siunitx}
\usepackage{hyperref}
\usepackage{float}
\usepackage[capitalize,noabbrev]{cleveref}

% --- Hyperref colors ---
\definecolor{linkblue}{RGB}{0,70,140}
\hypersetup{
  pdftitle={From Linear Scores to a Single Hidden Layer},
  pdfauthor={Sharaf Feyzullayev, Nijat Samadov, Samir Abdullazade},
  pdfsubject={Math4AI Capstone Report},
  pdfkeywords={softmax regression, neural networks, backpropagation,
               cross-entropy, classification, NumPy},
  colorlinks=true,
  linkcolor=linkblue,
  citecolor=linkblue,
  urlcolor=linkblue
}

% --- Better lists ---
\setlist[itemize]{itemsep=0.2em,topsep=0.3em}
\setlist[enumerate]{itemsep=0.2em,topsep=0.3em}
\sisetup{detect-all=true}

% ── Math shortcuts ──────────────────────────────────────────────────────────
\newcommand{\R}{\mathbb{R}}
\newcommand{\E}{\mathbb{E}}
\newcommand{\Prob}{\mathbb{P}}
\newcommand{\1}{\mathbf{1}}
\newcommand{\vone}{\mathbf{1}}
\newcommand{\T}{\mathsf{T}}
\newcommand{\loss}{\mathcal{L}}
\newcommand{\softmax}{\operatorname{softmax}}
\newcommand{\softmaxrow}{\operatorname{softmax}_{\mathrm{row}}}
\newcommand{\ce}{\operatorname{CE}}
\newcommand{\relu}{\operatorname{ReLU}}
\newcommand{\tanhact}{\operatorname{tanh}}

% ── Vectors and matrices ────────────────────────────────────────────────────
\newcommand{\vx}{\mathbf{x}}
\newcommand{\vy}{\mathbf{y}}
\newcommand{\vh}{\mathbf{h}}
\newcommand{\vs}{\mathbf{s}}
\newcommand{\vp}{\mathbf{p}}
\newcommand{\vb}{\mathbf{b}}
\newcommand{\vbone}{\mathbf{b}_1}
\newcommand{\vbtwo}{\mathbf{b}_2}
\newcommand{\vz}{\mathbf{z}}
\newcommand{\mX}{\mathbf{X}}
\newcommand{\mY}{\mathbf{Y}}
\newcommand{\mH}{\mathbf{H}}
\newcommand{\mP}{\mathbf{P}}
\newcommand{\mS}{\mathbf{S}}
\newcommand{\mW}{\mathbf{W}}
\newcommand{\mWone}{\mathbf{W}_1}
\newcommand{\mWtwo}{\mathbf{W}_2}
\newcommand{\mZ}{\mathbf{Z}}
\newcommand{\bdelta}{\boldsymbol{\delta}}
\newcommand{\bmu}{\boldsymbol{\mu}}

% ── Helpers ─────────────────────────────────────────────────────────────────
\newcommand{\etal}{\emph{et al.}}
\newcommand{\ie}{\emph{i.e.}, }
\newcommand{\eg}{\emph{e.g.}, }
\newcommand{\vsfull}{\emph{vs.}}

% ══════════════════════════════════════════════════════════════════════════════
%  TITLE / AUTHORS
% ══════════════════════════════════════════════════════════════════════════════

\title{\ProjectTitle}
\subtitle{\ProjectType}

\author{Sharaf Feyzullayev}
\affiliation{%
  \institution{\CourseName}
  \country{Baku, Azerbaijan}
}
\email{sheref.feyzullayev25@aiacademy.az}

\author{Nijat Samadov}
\affiliation{%
  \institution{\CourseName}
  \country{Baku, Azerbaijan}
}
\email{nicat.semedov25@aiacademy.az}

\author{Samir Abdullazade}
\affiliation{%
  \institution{\CourseName}
  \country{Baku, Azerbaijan}
}
\email{samir.abdullazade25@aiacademy.az}

% ══════════════════════════════════════════════════════════════════════════════
\begin{document}
% ══════════════════════════════════════════════════════════════════════════════

% ─── Abstract ────────────────────────────────────────────────────────────────
\begin{abstract}
We investigate when a one-hidden-layer neural network with $\tanh$ activation
genuinely improves upon multiclass softmax regression, and when additional model
complexity is unnecessary.  Both models are implemented from scratch using only
NumPy and evaluated on three tasks of increasing geometric difficulty: linearly
separable Gaussian blobs, nonlinear moons, and a 10-class digits benchmark.  On
the Gaussian task, softmax regression achieves identical accuracy (95.0\,\%) to
the neural network, confirming that nonlinearity adds no benefit when the data
is linearly separable.  On the moons task, the neural network achieves 97.5\,\%
accuracy versus the linear model's 83.8\,\%, demonstrating that the hidden
layer's nonlinear coordinate transform is essential for capturing curved decision
boundaries.  On digits, repeated-seed evaluation (5 seeds) shows the neural
network ($95.71\%\pm 0.91\%$) consistently outperforms softmax
($93.26\%\pm 0.73\%$).  A failure analysis with $H{=}2$ hidden units
(58.97\,\% accuracy) illustrates that capacity must match task complexity.
PCA/SVD analysis reveals that the digits' linear structure is largely captured
in ${\sim}20$ principal components, yet the neural network still provides a
meaningful representational advantage beyond what linear dimensionality
reduction achieves.
\end{abstract}

\keywords{softmax regression, neural networks, backpropagation, cross-entropy,
          classification, NumPy, PCA, capacity ablation}

\maketitle

% ══════════════════════════════════════════════════════════════════════════════
\section{Introduction}
% ══════════════════════════════════════════════════════════════════════════════

\paragraph{Central question.}
When does a one-hidden-layer nonlinear classifier genuinely improve on a linear
decision rule, and when is additional model complexity unnecessary?

This capstone investigates that question through a deliberately controlled
comparison between a linear probabilistic classifier and a simple nonlinear
neural model.  A responsible investigation requires (i)~deriving the mathematics
underlying each model class, (ii)~constructing experiments whose geometry
isolates the relevant differences, (iii)~evaluating with repeated seeds and
honest statistics, and (iv)~acknowledging where conclusions do and do not extend.
We study two model classes, each producing a categorical distribution over $K$
classes:
\begin{enumerate}[label=(\alph*),leftmargin=*]
  \item \textbf{Softmax Regression} — an affine map $\vz{=}\mW\vx{+}\vb$
        followed by softmax.  The decision boundary is a set of hyperplanes,
        so this model can only separate linearly distinguishable classes.
  \item \textbf{One-Hidden-Layer Neural Network} — a composition of an affine
        map, the $\tanh$ nonlinearity, and a second affine map followed by
        softmax.  The hidden layer performs a learned nonlinear coordinate
        change; the resulting decision boundary can curve, with complexity
        depending on the hidden width~$H$.
\end{enumerate}

\paragraph{Contributions.}
\begin{enumerate}[leftmargin=*]
  \item \textbf{Mathematical derivations from first principles.}
        We derive the softmax cross-entropy loss as negative log-likelihood
        (Section~\ref{sec:background}), derive the complete backpropagation
        equations for both model classes in vectorized mini-batch form, and
        provide a geometric argument explaining why each dataset matches or
        mismatches a linear decision rule~\cite{deisenroth2020mml, murphy2022pml}.

  \item \textbf{From-scratch NumPy implementation.}
        Both models and all three optimisers (SGD, Momentum, Adam) are
        implemented entirely in NumPy without any automatic-differentiation
        library.  The implementation covers model forward and backward passes,
        data loading and mini-batch generation, decision-boundary and
        loss-curve visualisation, optimiser updates, a mini-batch training
        loop with validation-checkpoint saving, automated sanity checks, and
        the full experiment orchestration including the PCA/SVD track, the
        repeated-seed evaluation, and the failure analysis.  This ensures that
        every gradient computation is understood at the level of the chain
        rule and verified numerically.

  \item \textbf{Controlled comparison across three geometric regimes.}
        The Gaussian, moons, and digits tasks are chosen to span linearly
        separable, nonlinearly separable, and moderately complex real-world
        geometry respectively, enabling principled, evidence-driven conclusions
        about when nonlinearity helps.

  \item \textbf{Required ablations, failure case, and repeated-seed statistics.}
        We conduct a capacity ablation on moons ($H\in\{2,8,32\}$), an optimiser
        study on digits (SGD, Momentum, Adam), a failure-case study
        ($H{=}2$ bottleneck on digits), and five-seed 95\,\% confidence
        intervals on the digits benchmark.

  \item \textbf{Track A: PCA/SVD and input geometry.}
        SVD of the digits training matrix reveals that ${\sim}20$ principal
        components capture 89.6\,\% of variance.  Softmax on 20-PC compressed
        inputs achieves 94.08\,\%, confirming that the neural network's
        advantage (${\sim}2.5$ ppt) originates from exploiting residual
        nonlinear structure.
\end{enumerate}

\paragraph{Preview of findings.}
The linear model is already sufficient on the Gaussian task (identical accuracy)
because the Bayes-optimal boundary is a hyperplane.  On moons, the hidden layer
provides an essential nonlinear coordinate change: softmax achieves 83.8\,\%
while the neural network reaches 97.5\,\%.  On digits, a statistically robust
${\sim}2.5$-point accuracy gain and ${\sim}42\%$ cross-entropy reduction favour
the neural network, with tight confidence intervals confirming the gap is not
a lucky-seed artifact.

% ══════════════════════════════════════════════════════════════════════════════
\section{Background and Mathematical Setup}
\label{sec:background}
% ══════════════════════════════════════════════════════════════════════════════

\subsection{Supervised Classification and Data Splits}

We assume labeled pairs $(\vx^{(i)}, y^{(i)})_{i=1}^{n}$ with $\vx^{(i)}\in\R^d$
and $y^{(i)}\in\{1,\dots,K\}$.  The training set is used for gradient updates;
the \emph{validation} set is used exclusively for selecting checkpoints and
hyperparameters; the \emph{test} set is evaluated exactly once per seed for
final reporting.  This three-way split prevents data leakage from model
selection into the reported performance.

\subsection{Softmax Regression}
\label{sec:softmax}

\subsubsection*{Forward pass}

Given $\vx\in\R^d$, softmax regression computes a logit vector:
\begin{equation}\label{eq:logits}
  \vz = \mW\vx + \vb, \qquad \mW\in\R^{K\times d},\;\vb\in\R^K.
\end{equation}
Converting logits to probabilities via softmax:
\begin{equation}\label{eq:softmax}
  \hat{y}_k = P(y{=}k\mid\vx;\mW,\vb)
  = \frac{\exp(z_k)}{\sum_{j=1}^{K}\exp(z_j)}, \quad k=1,\dots,K.
\end{equation}
Numerical stability is achieved by subtracting the row-wise maximum before
exponentiation, exploiting the shift-invariance
$\softmax(\vz)=\softmax(\vz-c\vone)$.  For a minibatch
$\mX\in\R^{n\times d}$:
\begin{align}
  \mS &= \mX\mW^\T + \vone\vb^\T \;\in\R^{n\times K}, \label{eq:sr_scores}\\
  \mP &= \softmaxrow(\mS)          \;\in\R^{n\times K}. \label{eq:sr_probs}
\end{align}

\subsubsection*{Cross-entropy as negative log-likelihood}

The cross-entropy loss with one-hot labels $\mY\in\{0,1\}^{n\times K}$ is:
\begin{equation}\label{eq:ce}
  J(\theta) = -\frac{1}{n}\sum_{i=1}^{n}\sum_{k=1}^{K}
              Y_{ik}\log P_{ik}.
\end{equation}
Minimising~\eqref{eq:ce} is algebraically identical to maximising the
conditional log-likelihood $\sum_{i}\log p_{y^{(i)}}(\vx^{(i)})$.  Adding
$L_2$ regularisation $\frac{\lambda}{2}\|\mW\|_F^2$ corresponds to MAP
estimation under a Gaussian prior~\cite{murphy2022pml, eth_intro_ml}; we use $\lambda=10^{-4}$ throughout.

\subsubsection*{Backward pass}

The combined softmax-plus-cross-entropy derivative simplifies to:
\begin{equation}\label{eq:sr_delta}
  \bdelta = \frac{\partial J}{\partial \mS}
  = \frac{1}{n}(\mP - \mY) \;\in\R^{n\times K}.
\end{equation}
Parameter gradients follow from the chain rule:
\begin{align}
  \frac{\partial J}{\partial \mW} &= \bdelta^\T\mX + \lambda\mW
    \;\in\R^{K\times d}, \label{eq:sr_dW}\\[3pt]
  \frac{\partial J}{\partial \vb} &= \bdelta^\T\vone
    \;\in\R^{K}. \label{eq:sr_db}
\end{align}

\subsection{One-Hidden-Layer Neural Network}
\label{sec:nn}

\subsubsection*{Forward pass}

The hidden layer applies an affine transformation followed by $\tanh$~\cite{cs231n_nn1, cs229dlnotes}:
\begin{align}
  \mZ_1    &= \mX\mWone^\T + \vone\vbone^\T \;\in\R^{n\times H},
             \label{eq:nn_Z1}\\
  \mH      &= \tanhact(\mZ_1)                \;\in\R^{n\times H}.
             \label{eq:nn_A1}
\end{align}
The output layer applies softmax regression to these transformed features:
\begin{align}
  \mS      &= \mH\mWtwo^\T + \vone\vbtwo^\T \;\in\R^{n\times K},
             \label{eq:nn_S}\\
  \mP      &= \softmaxrow(\mS)               \;\in\R^{n\times K}.
             \label{eq:nn_P}
\end{align}

\subsubsection*{Backpropagation}

\paragraph{Step 1: Output-layer error.}
\begin{equation}\label{eq:nn_delta2}
  \bdelta^{(2)} = \frac{1}{n}(\mP - \mY) \;\in\R^{n\times K}.
\end{equation}

\paragraph{Step 2: Output-layer gradients.}
\begin{align}
  \frac{\partial J}{\partial \mWtwo} &=
    (\bdelta^{(2)})^\T\mH + \lambda\mWtwo, \label{eq:nn_dW2}\\[3pt]
  \frac{\partial J}{\partial \vbtwo} &=
    (\bdelta^{(2)})^\T\vone. \label{eq:nn_db2}
\end{align}

\paragraph{Step 3: Propagate through the hidden layer.}
Error flows backwards through $\mWtwo$ and the local $\tanh$ derivative
$\tanhact'(z)=1-\tanhact^2(z)$:
\begin{equation}\label{eq:nn_delta1}
  \bdelta^{(1)} = \frac{\partial J}{\partial\mZ_1}
  = (\bdelta^{(2)}\mWtwo)\odot(1-\mH\odot\mH)
  \;\in\R^{n\times H}.
\end{equation}
Where $\tanh$ is saturated ($\mH\approx\pm 1$), the derivative is near zero,
producing the well-known \emph{vanishing gradient} effect~\cite{goodfellow2016deep, cs231n_backprop}.

\paragraph{Step 4: Hidden-layer gradients.}
\begin{align}
  \frac{\partial J}{\partial \mWone} &=
    (\bdelta^{(1)})^\T\mX + \lambda\mWone, \label{eq:nn_dW1}\\[3pt]
  \frac{\partial J}{\partial \vbone} &=
    (\bdelta^{(1)})^\T\vone. \label{eq:nn_db1}
\end{align}

\paragraph{Parameter count.}
Softmax regression has $K(d+1)$ parameters.  The neural network has
$H(d+1)+K(H+1)$.  For digits ($d{=}64$, $H{=}32$, $K{=}10$): 650 versus
2\,410 — roughly $3.7\times$ more capacity, which must be justified by
improved performance.

\subsection{Geometric Expectations for the Datasets}
\label{sec:geometry}

\paragraph{Why the Gaussian task matches a linear boundary.}
For two equal-covariance Gaussian classes
$\mathcal{N}(\bmu_0,\sigma^2\mathbf{I})$ and
$\mathcal{N}(\bmu_1,\sigma^2\mathbf{I})$, the log-posterior ratio is:
\begin{equation}\label{eq:logodds}
  \log\frac{P(y{=}1\mid\vx)}{P(y{=}0\mid\vx)}
  = \frac{1}{\sigma^2}(\bmu_1-\bmu_0)^\T\vx
  - \frac{1}{2\sigma^2}(\|\bmu_1\|^2-\|\bmu_0\|^2).
\end{equation}
This expression is \emph{affine} in $\vx$: the optimal boundary is a
hyperplane.  Softmax regression can represent this function exactly; a hidden
layer merely over-parameterises the same solution.

\paragraph{Why the moons task requires nonlinearity.}
The two classes form interleaving half-circles.  No single hyperplane can
separate them.  The hidden layer addresses this by learning:
\begin{equation}
  \vx \mapsto \tanhact(\mWone\vx+\vbone)\;\in\R^H,
\end{equation}
which ``untangles'' the arcs into linearly separable clusters in the
$H$-dimensional hidden space.  The quality of this untangling depends on
$H$, investigated in Section~\ref{sec:experiments}.

% ══════════════════════════════════════════════════════════════════════════════
\section{Methods}
\label{sec:methods}
% ══════════════════════════════════════════════════════════════════════════════

\subsection{Datasets}
\begin{itemize}[leftmargin=*]
  \item \textbf{Linear Gaussian} (\texttt{linear\_gaussian.npz}): 400 points
        in $\R^2$, two isotropic Gaussian blobs with mild overlap.
  \item \textbf{Moons} (\texttt{moons.npz}): 400 points in $\R^2$, two
        interleaving half-circles with additive noise.
  \item \textbf{Digits} (\texttt{digits\_data.npz},
        \texttt{digits\_split\_indices.npz}): $8\times 8$ pixel features,
        $d{=}64$, $K{=}10$, fixed split: 1\,258 train / 270 validation /
        269 test.
\end{itemize}

\subsection{Models}
Both models are initialised with Xavier weights ($W\sim\mathcal{N}(0,\sigma^2)$,
$\sigma{=}\sqrt{2/(d_\text{in}+d_\text{out})}$) and use the numerically stable
row-max-subtracted softmax.  The neural network uses hidden width $H{=}32$ as
the default for the digits benchmark and $H\in\{2,8,32\}$ for the moons
capacity ablation.

\subsection{Training Protocol}
Mini-batch SGD with batch size 64.  $L_2$ regularisation $\lambda{=}10^{-4}$.
Digits: 200-epoch budget; model selected by lowest \emph{validation}
cross-entropy checkpoint.  Synthetic tasks: 300 epochs, full-batch gradient
descent with $\lambda{=}10^{-5}$; the softmax baseline uses SGD
($\eta{=}0.1$) and the neural network uses Adam ($\eta{=}0.01$).
For the digits benchmark and repeated-seed study, the softmax baseline
uses SGD ($\eta{=}0.05$) and the neural network uses Adam ($\eta{=}0.001$),
selected as the best-calibrated optimiser from the optimiser study
(Section~\ref{sec:optimizer}).

\subsection{Evaluation Metrics}
We report (i)~\textbf{test accuracy}
$\frac{1}{n}\sum_{i}\mathbf{1}[\hat{y}^{(i)}{=}y^{(i)}]$ and
(ii)~\textbf{mean test cross-entropy}~\eqref{eq:ce}.  Both are reported
because accuracy ignores calibration whereas CE reflects the quality of the
full predicted distribution.

\subsection{Implementation Sanity Checks}
\label{sec:sanity}

Five automated checks verified code correctness before experimentation:
\begin{enumerate}[leftmargin=*]
  \item \textbf{Probability normalisation.}  Softmax rows sum to $1.0$ to
        machine precision on extreme logit vectors including
        $\vz{=}[-100,\,0,\,100]$.
  \item \textbf{Monotonic loss decrease.}  10 full-batch SGD steps on 50 samples
        produced strictly decreasing loss from 1.6202 to 1.3457.
  \item \textbf{Gradient verification.}  Analytical backpropagation matched
        centred finite differences ($\varepsilon{=}10^{-5}$) with a maximum
        discrepancy below $10^{-11}$ across all parameter matrices.
  \item \textbf{NaN/Inf detection.}  After 30 SGD steps of demanding training, all
        parameters remained finite.
  \item \textbf{Micro-dataset overfit.}  An $H{=}16$ network trained on
        10~examples for 200 Adam steps reached 100\,\% training accuracy,
        confirming sufficient expressivity.
\end{enumerate}

% ══════════════════════════════════════════════════════════════════════════════
\section{Experiments}
\label{sec:experiments}
% ══════════════════════════════════════════════════════════════════════════════

\subsection{Synthetic Decision Boundaries}

\paragraph{Linear Gaussian.}
Figure~\ref{fig:gaussian-boundaries} shows that both models learn essentially
the same straight boundary.  The neural network produces slightly different
probability contours but the 50\%-probability line is identical.  This confirms
Section~\ref{sec:geometry}: when the Bayes-optimal boundary is a hyperplane,
the hidden layer provides no accuracy advantage.

\begin{figure}[H]
  \centering
  \includegraphics[width=\columnwidth]{figures/comparison_lineargaussian.png}
  \caption{Decision boundaries on the Linear Gaussian task.  Left: Softmax
           Regression.  Right: Neural Network ($H{=}32$).  Both models produce
           a near-identical linear separator, confirming geometric
           compatibility.}
  \label{fig:gaussian-boundaries}
\end{figure}

\paragraph{Moons.}
Figure~\ref{fig:moons-boundaries} reveals the geometric difference.  Softmax
regression places a straight line through curved data, failing substantially.
The neural network maps $\R^2$ into $\R^{32}$ via $\tanhact(\mWone\vx+\vbone)$,
where the arcs become linearly separable, achieving near-100\,\% accuracy.

\begin{figure}[H]
  \centering
  \includegraphics[width=\columnwidth]{figures/comparison_moons.png}
  \caption{Decision boundaries on the Moons task.  Left: Softmax Regression
           (fails, 83.8\,\%).  Right: Neural Network ($H{=}32$, 97.5\,\%).
           The hidden layer's nonlinear transform ``untwists'' the crescent
           structure.}
  \label{fig:moons-boundaries}
\end{figure}

\subsection{Capacity Ablation on Moons}
\label{sec:ablation}

Figure~\ref{fig:ablation} shows the effect of hidden width on boundary quality:
\begin{itemize}[leftmargin=*]
  \item $H{=}2$: angular, low-confidence boundaries — insufficient dimensions.
  \item $H{=}8$: reasonable curved boundary, uneven confidence near edges.
  \item $H{=}32$: smooth, high-confidence boundary tightly tracking the data
        manifold.
\end{itemize}
This demonstrates a monotonic relationship between hidden width and boundary
quality \emph{on this task}.  On the Gaussian task, all three widths produce
essentially the same linear boundary because no additional representational
power is needed.

\begin{figure}[H]
  \centering
  \includegraphics[width=\columnwidth]{figures/capacity_ablation_boundaries.png}
  \caption{Capacity ablation on Moons.  Left to right: $H{=}2$, $H{=}8$,
           $H{=}32$.  Boundary smoothness and confidence improve monotonically
           with hidden width.}
  \label{fig:ablation}
\end{figure}

\subsection{Digits Benchmark}
\label{sec:digits}

Table~\ref{tab:digits_main} summarises the core comparison with five-seed
repeated evaluation.

\begin{table}[H]
  \centering
  \caption{Digits benchmark: five-seed summary (mean $\pm$ 95\,\% CI).}
  \label{tab:digits_main}
  \begin{tabular}{@{}lcc@{}}
    \toprule
    \textbf{Model} & \textbf{Test Accuracy} & \textbf{Test CE}\\
    \midrule
    Softmax Regression       & $93.26\%\pm 0.73\%$ & $0.2692\pm 0.0067$\\
    Neural Network ($H{=}32$)& $95.71\%\pm 0.91\%$ & $0.1566\pm 0.0162$\\
    \bottomrule
  \end{tabular}
\end{table}

The neural network improves accuracy by ${\sim}2.5$ percentage points and
reduces cross-entropy by ${\sim}42\%$.  The 95\,\% confidence intervals do not
overlap.  The CE reduction is proportionally larger than the accuracy gain,
indicating better-calibrated probabilities.  Figure~\ref{fig:loss_curves_digits}
shows representative training dynamics; neither model exhibits severe
overfitting within the 200-epoch budget.

\begin{figure}[H]
  \centering
  \includegraphics[width=\columnwidth]{figures/loss_curves_digits.png}
  \caption{Training and validation loss dynamics on the Digits benchmark
           (single representative seed; Softmax uses SGD $\eta{=}0.05$,
           NN uses Adam $\eta{=}0.001$).  Neither model shows severe
           overfitting within 200 epochs.}
  \label{fig:loss_curves_digits}
\end{figure}

\subsection{Optimizer Study on Digits}
\label{sec:optimizer}

Three optimisers were compared on the neural network ($H{=}32$):

\begin{table}[H]
  \centering
  \caption{Optimiser study on Digits (NN, $H{=}32$, best val-CE checkpoint).}
  \label{tab:optimizers}
  \begin{tabular}{@{}lccc@{}}
    \toprule
    \textbf{Optimiser} & \textbf{Test Acc.} & \textbf{Test CE} &
    \textbf{Best Epoch}\\
    \midrule
    SGD ($\eta{=}0.05$)               & 94.84\% & 0.1700 & 197\\
    Momentum ($\eta{=}0.05,\mu{=}0.9$)& 95.38\% & 0.1653 & 197\\
    Adam ($\eta{=}0.001$)             & 95.11\% & 0.1507 & 197\\
    \bottomrule
  \end{tabular}
\end{table}

Adam and Momentum converge to low validation loss faster than plain SGD, but
all three reach comparable final accuracy.  Adam achieves the lowest CE
(0.1507), suggesting better calibration, while Momentum yields the highest
accuracy (95.38\,\%).  The practical conclusion: optimiser choice has a
secondary effect (${\sim}0.5$ ppt spread) relative to model-class choice
(${\sim}2.5$ ppt gap).

\begin{figure}[H]
  \centering
  \includegraphics[width=\columnwidth]{figures/optimizer_study_digits.png}
  \caption{Validation loss and accuracy across optimisers (NN, $H{=}32$,
           Digits).  Adam and Momentum converge faster; all three reach similar
           final accuracy.}
  \label{fig:optimizers}
\end{figure}

\subsection{Failure-Case Analysis: Under-Capacity Bottleneck ($H{=}2$)}
\label{sec:failure}

We deliberately trained a neural network with $H{=}2$ on the 10-class Digits
task.  The hidden layer projects the 64-dimensional input to a 2-dimensional
representation $\vh{=}\tanhact(\mWone\vx+\vbone)\in[-1,1]^2$ before the
output layer attempts 10-class classification.

\begin{table}[H]
  \centering
  \caption{Failure case: $H{=}2$ bottleneck on Digits.}
  \label{tab:failure}
  \begin{tabular}{@{}lcc@{}}
    \toprule
    \textbf{Model} & \textbf{Test Accuracy} & \textbf{Test CE}\\
    \midrule
    Softmax Regression (linear) & 93.26\% & 0.2692\\
    NN ($H{=}2$)                & 58.97\% & 1.1133\\
    NN ($H{=}32$)               & 95.71\% & 0.1566\\
    \bottomrule
  \end{tabular}
\end{table}

The $H{=}2$ model achieves only 58.97\,\%---far below even the linear baseline.
A nonlinear model \emph{worse} than the linear one.

\paragraph{Why it fails: a geometric argument.}
The failure is \emph{representational}, not optimisational.  With $H{=}2$, the
network maps $\R^{64}$ to $[-1,1]^2$.  From only two dimensions, the output
layer must carve 10 linearly separable regions---geometrically impossible for
the complexity of handwritten digits.  The optimiser faithfully minimises the
loss \emph{within} the representational constraint, but the global minimum of
the $H{=}2$ loss surface itself corresponds to a poor classifier.

\paragraph{Lesson.}
\textbf{Representational capacity must match task complexity before optimisation
can be effective.}  Increasing $H$ from 2 to 32 does not change the
optimiser---it changes the \emph{geometry of the function class}, enabling
richer internal representations.  An insufficient bottleneck can actually
\emph{destroy} information present in the original input, making the nonlinear
model worse than the linear baseline that preserves all $d{=}64$ dimensions.

\begin{figure}[H]
  \centering
  \includegraphics[width=\columnwidth]{figures/failure_analysis.png}
  \caption{Failure case.  The $H{=}2$ bottleneck (red) severely underfits with
           high residual loss, while the $H{=}32$ model (green) converges to
           low loss.  The gap is due to representational capacity, not
           optimisation quality.}
  \label{fig:failure}
\end{figure}

% ══════════════════════════════════════════════════════════════════════════════
\section{Advanced Track A: PCA/SVD and Input Geometry}
\label{sec:pca}
% ══════════════════════════════════════════════════════════════════════════════

We computed the SVD of the mean-centred training matrix
$\mX_c{=}\mX{-}\vone\bar{\vx}^\T{=}\mathbf{U}\boldsymbol{\Sigma}\mathbf{V}^\T$
to study linear structure before any learning occurs.

\paragraph{Scree analysis.}
Rapid singular-value decay occurs (Figure~\ref{fig:pca}a):
\begin{center}
\begin{tabular}{@{}lcccc@{}}
  \toprule
  Components $m$ & 5 & 10 & 20 & 40\\
  \midrule
  Cumulative variance & 54.7\% & 74.1\% & 89.6\% & 98.9\%\\
  \bottomrule
\end{tabular}
\end{center}
The top 20 components capture 89.6\,\% of total variance; the remaining 44
dimensions contribute noise-level structure.

\paragraph{2-D projection.}
Figure~\ref{fig:pca}b shows digits projected onto the first two PCs.  Digits 0
and 6 form distinct clusters; digits 3, 5, and 8 overlap (shared curved
strokes).  Two dimensions are insufficient for full separation but reveal
meaningful coarse structure.

\paragraph{Classification at reduced dimensions.}
Table~\ref{tab:pca} shows softmax accuracy on PCA-compressed inputs:

\begin{table}[H]
  \centering
  \caption{Softmax Regression accuracy on PCA-reduced Digits inputs.}
  \label{tab:pca}
  \begin{tabular}{@{}lcccc@{}}
    \toprule
    $m$ & 10 & 20 & 40 & 64 (full)\\
    \midrule
    Val.\ Accuracy & 92.68\% & 94.08\% & 94.93\% & 95.49\%\\
    Val.\ CE Loss  & 0.2874  & 0.2292  & 0.2133  & 0.2109\\
    \bottomrule
  \end{tabular}
\end{table}

At $m{=}20$, softmax still achieves 94.08\,\%---only 1.4 points below the
full-dimension result.  This reveals \textbf{substantial linear structure} in
the input space.  The neural network's ${\sim}2.5$-point advantage comes from
exploiting \emph{residual nonlinear} interactions that PCA---a purely linear
projection---cannot recover.

\begin{figure}[H]
  \centering
  \begin{subfigure}[b]{0.48\columnwidth}
    \centering
    \includegraphics[width=\textwidth]{figures/pca_scree_digits.png}
    \caption{Scree plot and cumulative explained variance.}
    \label{fig:pca_scree}
  \end{subfigure}\hfill
  \begin{subfigure}[b]{0.48\columnwidth}
    \centering
    \includegraphics[width=\textwidth]{figures/pca_2d_digits.png}
    \caption{2-D PCA projection of Digits classes.}
    \label{fig:pca_2d}
  \end{subfigure}
  \caption{Track A: PCA/SVD analysis of the Digits dataset.}
  \label{fig:pca}
\end{figure}

% ══════════════════════════════════════════════════════════════════════════════
\section{Discussion}
\label{sec:discussion}
% ══════════════════════════════════════════════════════════════════════════════

We answer the five core interpretive questions, drawing exclusively on the
evidence above.

\begin{enumerate}[leftmargin=*]

\item \textbf{When is the linear model geometrically sufficient, and when is it
not?}

The linear model suffices when the Bayes-optimal boundary is a hyperplane.
This holds rigorously for the Gaussian task (Equation~\ref{eq:logodds}) and
approximately for Digits, where softmax reaches 93.26\,\% and PCA confirms
that dominant discriminative directions are linear.  The linear model is
\emph{insufficient} on Moons, where no affine function separates the arcs
(Figure~\ref{fig:moons-boundaries}).

\item \textbf{What representational change does the hidden layer provide?}

The map $\vh{=}\tanhact(\mWone\vx{+}\vbone)$ is a learned nonlinear coordinate
change from $\R^d$ to $\R^H$.  On Moons, it ``untwists'' the arcs into linearly
separable clusters in $\R^{32}$.  On Digits, it captures nonlinear feature
combinations (stroke curvatures, junctions) that linear projections cannot,
accounting for the ${\sim}2.5$-point gain.  The PCA comparison confirms this
gain is from \emph{nonlinear} structure.

\item \textbf{When does additional model complexity fail to justify itself?}

\emph{(a)} When the baseline already captures the geometry: on the Gaussian
task, the NN matches softmax with $3.7\times$ more parameters.
\emph{(b)} When complexity is of the wrong kind: softmax at $m{=}20$ PCA
dimensions (94.08\,\%) is within 1.6 points of the full NN (95.71\,\%).
If engineering cost is high and the gap is tolerable, the simpler model may
be preferable.

\item \textbf{What does the failure case teach about optimisation, capacity, or
generalisation?}

The $H{=}2$ bottleneck teaches that \emph{capacity is a prerequisite that
optimisation cannot substitute}.  Identical optimiser settings yield 58.97\,\%
($H{=}2$) versus 95.71\,\% ($H{=}32$).  This distinguishes three failure modes:
\begin{itemize}[leftmargin=*]
  \item \emph{Optimisation failure}: architecture can represent a solution,
        optimiser cannot find it.
  \item \emph{Capacity failure}: architecture \emph{cannot} represent a
        solution. $\leftarrow$ this case.
  \item \emph{Generalisation failure}: model fits training but not test data.
\end{itemize}
The $H{=}2$ case is pure capacity failure: increasing $H$ expands the
hypothesis class, not the optimiser's effectiveness.

\item \textbf{What do repeated-seed statistics add beyond a single run?}

A single run depends on random initialisation and mini-batch order.  Five seeds
with $\bar{x}\pm 2.776\cdot s/\!\sqrt{5}$ bound \emph{expected} performance.
Tight intervals (softmax: $\pm 0.73\%$; NN: $\pm 0.91\%$) confirm the
${\sim}2.5$-point gap is robust, not a lucky-seed artifact.  The NN's wider
interval reflects greater initialisation sensitivity from additional parameters.

\end{enumerate}

% ══════════════════════════════════════════════════════════════════════════════
\section{Limitations}
\label{sec:limitations}
% ══════════════════════════════════════════════════════════════════════════════

\begin{enumerate}[leftmargin=*]
  \item \textbf{Architectural scope.}  Single hidden layer, $\tanh$ only,
        $H\leq 32$.  Conclusions do not extend to deeper networks, wider
        layers, or alternative activations (ReLU, GELU).

  \item \textbf{Dataset scale and splits.}  Digits has $d{=}64$, $K{=}10$,
        ${\sim}1\,800$ samples with a single fixed split.  Results may not
        generalise to high-resolution images, larger class counts, or
        cross-validated protocols.

  \item \textbf{Optimiser and regularisation exploration.}  Each optimiser was
        tested at one fixed learning rate; only $L_2$ ($\lambda{=}10^{-4}$)
        regularisation was used.  Different schedules, dropout, or batch
        normalisation could change relative rankings.

  \item \textbf{Repeated-seed sample size.}  Five seeds assume approximate
        normality of the accuracy distribution, which is difficult to verify
        with $n{=}5$.

  \item \textbf{Synthetic task simplicity.}  400 points, 2-D, noise-free
        labels---illustrative of geometric principles but not representative
        of real-world data ambiguity.  No formal calibration analysis
        (reliability diagrams, ECE) was performed.
\end{enumerate}

Our findings reliably characterise the specific comparison on these three tasks
under this fixed protocol.  They should not be overgeneralised to deeper
architectures, larger datasets, or more sophisticated training pipelines.

% ══════════════════════════════════════════════════════════════════════════════
\vspace{-0.3cm}
\section*{Acknowledgments}
We thank the instructors and reviewers of the Math4AI capstone for guidance
and feedback throughout the project.
\vspace{-0.5cm}
\setlength{\bibsep}{5pt}
\nocite{*}
\bibliographystyle{ACM-Reference-Format}
\bibliography{references}

% ══════════════════════════════════════════════════════════════════════════════
\appendix
% ══════════════════════════════════════════════════════════════════════════════
\vspace{-0.3cm}
\section{Supplementary Figures}
\label{app:figures}

\subsection{Linear Gaussian Decision Boundaries}

The following plots illustrate the decision boundaries learned by Softmax
Regression and the Neural Network on the Linear Gaussian dataset. Since
the data is purely linearly separable, both models correctly converge to
a straight-line boundary, providing near-identical classification
performance and showing that additional non-linear capacity offers no
structural advantage here.

\begin{figure}[H]
  \centering
  \includegraphics[width=0.65\columnwidth]{figures/decision_boundary_lineargaussian_softmax.png}
  \caption{Softmax Regression decision boundary on the Linear Gaussian task.}
  \label{app:fig:gaussian_softmax}
\end{figure}

\begin{figure}[H]
  \centering
  \includegraphics[width=0.65\columnwidth]{figures/decision_boundary_lineargaussian_nn.png}
  \caption{Neural Network ($H{=}32$) decision boundary on the Linear Gaussian task.}
  \label{app:fig:gaussian_nn}
\end{figure}

\subsection{Moons Decision Boundaries}

Conversely, the Moons dataset poses a non-linear classification problem. The
linear constraints of Softmax Regression force it to draw a suboptimal line
across intertwined classes. In contrast, the Neural Network leverages its
hidden layer to non-linearly transform the input space, effectively wrapping
the decision boundary around the clusters.

\begin{figure}[H]
  \centering
  \includegraphics[width=0.65\columnwidth]{figures/decision_boundary_moons_softmax.png}
  \caption{Softmax Regression decision boundary on the Moons task.}
  \label{app:fig:moons_softmax}
\end{figure}

\begin{figure}[H]
  \centering
  \includegraphics[width=0.65\columnwidth]{figures/decision_boundary_moons_nn.png}
  \caption{Neural Network ($H{=}32$) decision boundary on the Moons task.}
  \label{app:fig:moons_nn}
\end{figure}

\subsection{Moons Capacity Ablation}

We trained the Neural Network on the Moons task varying the hidden units
($H\in\{2,8,32\}$). A narrow bottleneck ($H{=}2$) struggles to capture
continuous curvature, resulting in sharp boundaries. Increasing the width
monotonically improves both boundary smoothness and classification certainty.

\begin{figure}[H]
  \centering
  \begin{subfigure}[b]{0.65\columnwidth}
    \centering
    \includegraphics[width=\textwidth]{figures/decision_boundary_moons_h2.png}
    \caption{$H{=}2$: insufficient capacity, angular boundary.}
  \end{subfigure}

  \vspace{0.6em}

  \begin{subfigure}[b]{0.65\columnwidth}
    \centering
    \includegraphics[width=\textwidth]{figures/decision_boundary_moons_h8.png}
    \caption{$H{=}8$: reasonable curved boundary.}
  \end{subfigure}

  \vspace{0.6em}

  \begin{subfigure}[b]{0.65\columnwidth}
    \centering
    \includegraphics[width=\textwidth]{figures/decision_boundary_moons_h32.png}
    \caption{$H{=}32$: smooth, high-confidence boundary.}
  \end{subfigure}
  \caption{Capacity ablation on Moons — individual panels.}
  \label{app:fig:ablation}
\end{figure}

\subsection{Loss Dynamics}

These learning curves display the cross-entropy loss over time for the
synthetic datasets. We observe stable monolithic decreases in both training
and validation losses, confirming effective gradient descent without
significant overfitting within the allocated epochs.

\begin{figure}[H]
  \centering
  \includegraphics[width=0.65\columnwidth]{figures/loss_curves_lineargaussian.png}
  \caption{Training and validation loss dynamics for the Linear Gaussian task.}
  \label{app:fig:lc_gaussian}
\end{figure}

\begin{figure}[H]
  \centering
  \includegraphics[width=0.65\columnwidth]{figures/loss_curves_moons.png}
  \caption{Training and validation loss dynamics for the Moons task.}
  \label{app:fig:lc_moons}
\end{figure}

\begin{figure}[H]
  \centering
  \includegraphics[width=0.65\columnwidth]{figures/capacity_ablation_curves.png}
  \caption{Training curves for the capacity ablation ($H\in\{2,8,32\}$) on Moons.}
  \label{app:fig:ablation_curves}
\end{figure}

\subsection{Repeated-Seed and PCA Analysis}

Statistical confidence metrics over repeated runs confirm the strict separation
between the models' accuracies. PCA validates that while dominant characteristics
project well linearly, residual non-linear dependencies drive the performance gap.

\begin{figure}[H]
  \centering
  \includegraphics[width=0.65\columnwidth]{figures/repeated_seed_digits.png}
  \caption{Repeated-seed evaluation on Digits (5 seeds, 95\,\% CI).}
  \label{app:fig:repeated_seed}
\end{figure}

\begin{figure}[H]
  \centering
  \includegraphics[width=0.65\columnwidth]{figures/pca_softmax_comparison.png}
  \caption{Softmax accuracy as a function of PCA dimensions.}
  \label{app:fig:pca_softmax}
\end{figure}

% ══════════════════════════════════════════════════════════════════════════════
\section{Contribution Statement}
\label{app:contributions}
% ══════════════════════════════════════════════════════════════════════════════

\begin{itemize}[leftmargin=*]
  \item \textbf{Sharaf Feyzullayev} implemented the two core model classes
        (Softmax Regression and One-Hidden-Layer Neural Network), the
        data-loading and mini-batch generation utilities, and the full
        decision-boundary and loss-curve plotting functions.  Sharaf
        also led the mathematical derivations for forward and backward passes
        and designed the moons capacity-ablation experiment.

  \item \textbf{Nijat Samadov} implemented the full experiment orchestration
        script, covering both synthetic tasks (Linear Gaussian evaluation)
        and the advanced Track A (PCA/SVD scree plot, 2-D projection,
        reduced-dimension softmax comparison) as well as the repeated-seed
        evaluation with 95\,\% confidence intervals and the $H{=}2$ failure
        analysis.

  \item \textbf{Samir Abdullazade} implemented the three optimisers (SGD,
        Momentum, Adam) from scratch, wrote the mini-batch training loop
        with validation-checkpoint saving, and designed the five automated
        sanity checks.  Samir also ran the digits optimiser comparison study
        and produced the associated figures and result logs.

\end{itemize}

\end{document}
